\section{The first split one-triple application}\label{sec:first-c1}
We consider subgroups with exactly one actual regular $C_3$ orbit and
no actual natural $A_4$ orbit. On its complement of size $m=n-3$, an
external graph is determined by a literal nonzero character
$\chi:K\to C_3$. Fix a generator of the displayed $C_3$ to specify
its orientation. For split characters, let $\Theta(K)$ be the set
whose complete graphs survive any preceding disjoint exclusions.
The exact surviving character weight is
\begin{equation}\label{eq:c1-surviving-weight}
 \sigma^{\mathrm{surv}}_3(K)=|\Theta(K)|
 =\sum_{\substack{\chi\ne0\\\chi\text{ split}}}
       {\bf1}_{\{\operatorname{graph}(\chi)\text{ survives}\}},
 \qquad
 0\leq\sigma^{\mathrm{surv}}_3(K)\leq3^{d_3(K)}-3^{h_3(K)}.
\end{equation}
In particular the inequality need not be equality after exclusions.
On a fixed labelled complement the surviving mass is precisely
\begin{equation}\label{eq:c1-exact-mass}
 \frac{n!}{6m!}\sum_{K\leq S_m}\sigma^{\mathrm{surv}}_3(K).
\end{equation}
The sum restricts to residuals without a regular $C_3$ or natural
$A_4$ orbit. The factor six is the original action normalizer, and
there is one selected triple because the actual orbit is unique.

\begin{proposition}[Low one-triple cone]\label{prop:c1-low}
The split graphs with $20d_3(K)\leq3m$ have normalized mass at most
\[
 K_{\mathrm{lo}}(n,m)\frac{s_m}{L_m},\qquad
 K_{\mathrm{lo}}(n,m)=\frac16\Delta(n,m)3^{3m/20}
 \leq2^{-\eta m+O(\log(n+2))},
 \quad \eta=\frac14-\frac{3\log_2 3}{20}>0.
\]
\end{proposition}
\begin{proof}
Apply \eqref{eq:c1-surviving-weight} to each residual before enlarging
its number to $s_m$. The labelled factorials cancel exactly as in
\eqref{eq:c1-exact-mass}. Here is an explicit bound covering every
fixed deletion width $w$:
\begin{equation}\label{eq:arbitrary-width-pointing}
 \Delta(b+w,b)\leq
 \phi^{-2}2^{(w+1)/4}(b+w+1)^{w+2}
 2^{-\lfloor w/2\rfloor b/4},
 \qquad \phi=\prod_{j\geq1}(1-2^{-j}).
\end{equation}
Indeed, the two guarded parity coefficients at rank $k$ both lie
between $A_k$ and $(k+1)A_k$, since $A_{k-1}\leq2kA_k$.
The rank shift $t=\lfloor(b+w)/2\rfloor-\lfloor b/2\rfloor$
lies between $\lfloor w/2\rfloor$ and $w$. Apply the coefficient
shift bound $A_k\leq(2(k+t))^tA_{k+t}$ and the Gaussian ratio bound;
their polynomial factors are at most $(b+w+1)^{w+2}$.
The Gaussian exponent is at most
$(w+1)/4-\lfloor w/2\rfloor b/4$, proving the display also at rank zero.
For $w=3$ this gives the required loss $m/4$. Finally $3^5<2^8$
gives $3\log_2(3)/20<6/25<1/4$, so the low-cone kernel is also
$O(2^{-m/200})$ after absorbing its fixed polynomial factor.
\end{proof}

We next prove the earlier orbit estimates actually needed for the high
cone. Only finitely many physical degrees occur in this application;
no growing action menu is needed. Put $\gamma=(\log_2 3)/3$ and
\[
 Z_J(U)=\sum_{N\lhd U}|\operatorname{Epi}(J,U/N)|,
 \qquad a(U)=|N_{S_w}(U)|,\quad J\leq S_b.
\]
All target normal subgroups in this sum are literal. Fixing each target
quotient map is the asymmetric Goursat convention, so there is no
additional automorphism divisor.

\begin{lemma}[Finite earlier orbit fibres]\label{lem:c1-earlier-fibres}
For the following actual actions, $Z_J(U)/a(U)$ is at most
$D_U(1+b)2^{\alpha_U b}$, with a fixed constant $D_U$:
\[
\begin{array}{c|c|c}
 \text{action}&w&\alpha_U\\\hline
 6T1,6T4,6T5,6T6&6&\gamma\\
 12T20,12T85,12T164,12T194&12&2\gamma\\
 12T228,12T229,12T265&12&1/2+5\log_2(3)/9\\
 \text{transitive }3\text{-group}&9&5\log_2(3)/9\\
 \text{transitive }3\text{-group}&27&5\log_2(3)/3
\end{array}
\]
The degree-nine row and the three final degree-twelve actions require
$J$ to have exactly one regular $C_3$ orbit and no natural $A_4$
orbit. Every other row permits arbitrary $J$.
\end{lemma}
\begin{proof}
We give the fibre arguments, including the shared source budgets.
We use the permutation prime-rank bound $d_p(F)\leq b/p$ for every
literal subgroup $F\leq S_b$, and the abelianization-order bound
$|J/J'|\leq3^{b/3}$ from \cite{KovacsPraeger1989}. Constants depending
on a fixed target, including its finite normal menu, are absorbed
only after its full quotient sum has been written.

For $6T1=C_6$, the full fibre is $|\operatorname{Hom}(J,C_6)|$.
For $6T5=C_3\wr C_2\cong C_3\times S_3$, its normal quotients
are $1,C_2,C_3,C_6,S_3,C_3\times S_3$. Set $F=J'J^2$ and
$E=J/F$. The semisimple $\mathbb F_3E$-module $H_1(F;\mathbb F_3)$
has trivial multiplicity $d=d_3(J)$ and sign multiplicities
$a_\theta$, with $d+\sum a_\theta\leq b/3$. Coprime
inflation--restriction gives exactly
\[
 |\operatorname{Epi}(J,S_3)|=3\sum_{\theta\ne0}(3^{a_\theta}-1),
\quad
 |\operatorname{Epi}(J,C_3\times S_3)|
 =(3^d-1)3\sum_{\theta\ne0}(3^{a_\theta}-1).
\]
The product is onto because $C_3$ and $S_3$ have no common
nontrivial quotient. Since $\sum(3^{a_\theta}-1)\leq
3^{\sum a_\theta}-1$, each expression is at most $3^{b/3+1}$.
The central and twisted characters have used one module budget.

The actions $6T4,6T6$ are respectively $A_4$ in degree six and
$C_2\times A_4$ in degree six. Their binary bases are cyclic
$\mathbb F_2[C_3]$-modules. Likewise $12T20,12T85,12T164$ have
cyclic binary bases and elementary tops $C_3^2$, with the top
action on the base allowed to have a kernel. Here is the requisite
general estimate for $M\rtimes C_3^r$, $r=1,2$, with cyclic binary
$M$. Each quotient is $V\rtimes C_3^h$, $h\leq r$, with $V$
a multiplicity-one semisimple module. Use the same source
$F=J'J^3$, $E=J/F$, $D=H_1(F;\mathbb F_2)$ for every top map.
Write the trivial multiplicity of $D$ as $r_0$ and its nontrivial
simple multiplicities as $a_\alpha$. Each nontrivial type has
dimension two, and
\[
 r_0+2\sum_\alpha a_\alpha\leq b/2.
\]
For a fixed onto top map $\theta:E\to C_3^h$, surjective
restrictions to $F$ number
$(2^{r_0}-1)^e\prod_i(4^{a_{\alpha_i(\theta)}}-1)$, where
$e\in\{0,1\}$. If a nontrivial target constituent is retained,
fix one of its two concrete ternary characters. Pulling it back
to a specified source type permits at most
$2\cdot3^{(h-1)d_3(J)}$ top maps. Bound the remaining factors
using the one source budget and sum the retained type. The result
is at most
\[
 (b/2)\,3^{(h-1)d_3(J)}2^{b/2}.
\]
For each restriction and top map, the lift torsor has at most
$|M|$ elements, by coprime cohomology on $E$. Abelian quotients
are bounded by $|J/J'|^r$ (or by one such factor when $h=0$).
Thus all quotients obey the claimed $(1+b)2^{r\gamma b}$ bound.
This argument treats nonfaithful top actions as well as faithful ones.

For $12T194=C_3\wr A_4$, use its full ternary block kernel.
Every quotient is an extension of a quotient $S$ of $A_4$ by a
module subquotient $X$ of the natural permutation module
$\mathbb F_3^4$. Every such $X$ embeds in $\mathbb F_3[S]$.
To see this, first restrict to the normal $V_4$. Over a splitting
field its trivial character gives one trivial $C_3$-module;
its three nontrivial characters form one orbit and give a simple
matrix-algebra component. Subquotients of the first embed in the
regular $C_3$ block, and the second component is semisimple.
The whole permutation module embeds in $\mathbb F_3[A_4]$ by
coset sums. If the action factors through $S$, the image lies in
the fixed space of the kernel, which is the inflated regular
$\mathbb F_3[S]$-module. For a concrete top map $J\to S$ with
kernel $F$, the restriction choices are therefore at most
$3^{d_3(F)}\leq3^{b/3}$. The compatible lift torsor has a fixed
target-dependent size: after the restriction is fixed, differences
descend to $S$. The preceding cyclic-module estimate bounds the
top maps by $D2^{\gamma b}$. This gives $D2^{2\gamma b}$.

For degree-nine $3$-groups, embed $U$ in $C_3\wr C_3$ and put
$V=U\cap C_3^3$. It is a submodule of the uniserial ring
$\mathbb F_3[t]/(t^3)$, so $d_3(U)\leq2$. Each quotient is
elementary ternary of rank at most two, or is an extension of
$C_3$ by a module $M$ embedding in its regular module.
For a nonzero $\theta:J\to C_3$, Shapiro's lemma gives
\[
 |Z^1(J,\mathbb F_3[C_3]_\theta)|
   =3^{d_3(\ker\theta)+2}.
\]
The actual lift set is empty or a torsor for a subspace of this
cocycle space; this remains valid for a nonsplit target. Thus every
quotient $Q$ satisfies
\begin{equation}\label{eq:degree-nine-joint}
 |\operatorname{Epi}(J,Q)|\leq3^{d_3(J)+b/3+2}.
\end{equation}
By \eqref{eq:rank3-pattern}, the stipulated pattern gives
$d_3(J)\leq(2b+3)/9$, proving the degree-nine row. The more exact
quantity before compression is the single sum
$\sum_{\theta\ne0}3^{d_3(\ker\theta)+2}$.

For the three remaining degree-twelve actions, the actual action has
three blocks of four points, binary translation base $W\leq V_4^3$,
and a $3$-group top $P$ acting transitively on the nine nonzero local
translations. The full module $V_4^3$ is cyclic: a nonzero local
translation generates its local plane under the block stabilizer,
and its top translates generate all three planes. Every invariant
submodule is cyclic by Maschke's theorem. Consequently, for every
normal quotient $Q$, its binary kernel $M$ is cyclic for its top
quotient $S$ of $P$. Put $F=O^3(J)$ and $E=J/F$. For a fixed
map $J\to S$, the pulled-back $M$ embeds in the regular
$\mathbb F_2[E]$-module. It follows that
\[
 |\operatorname{Hom}_E(H_1(F;\mathbb F_2),M)|\leq2^{b/2}.
\]
Coprime lift torsors on $E$ cost at most $|M|\leq2^6$.
Apply \eqref{eq:degree-nine-joint} to the same source $J$ and the
actual nine-point action of $P$, before enlarging the pattern class
of sources. This gives the exponent $1/2+5\log_2(3)/9$.

Finally we prove the degree-27 row without any classification of
its $3$-groups. More generally, for a transitive $3$-group of degree
$w=3^k$, every quotient satisfies
\begin{equation}\label{eq:ternary-quotient-small}
 |\operatorname{Epi}(J,Q)|\leq D_w3^{(w+3)b/18},
\end{equation}
where $D_w$ is independent of $J$ and $U$. Choose a central
semiregular $C_3$ in $U$; its blocks have full elementary ternary
kernel $V\leq\mathbb F_3^s$, $s=w/3$, and transitive $3$-group
top $T$. A central semiregular $C_3$ of $T$ makes this permutation
module free of rank $s/3$ on restriction. Every subquotient $M$
therefore has socle dimension at most $s/3$. If its action factors
through $S$, it embeds in at most $s/3$ copies of the regular
$\mathbb F_3[S]$-module: the group algebra is a self-injective
algebra with its unique simple module trivial. Hence, over a fixed
map $J\to S$, restrictions from its literal kernel number at most
$3^{sb/9}$. After fixing a restriction, lift differences descend
to $S$ and have a bounded target-dependent cost. Induction on $k$
gives \eqref{eq:ternary-quotient-small}, since
$(s+3)/18+s/9=(w+3)/18$. If $S=1$, use the simultaneous head
bound $d_3(U)\leq(w+3)/6$, obtained from the same block filtration.
The base $w=3$ is the elementary rank bound. For $w\leq27$ the
normal menus and target sizes are finite, so their uniform constants
are finite. At $w=27$ the coefficient is $5/3$, as required.
\end{proof}

\begin{lemma}[Actual action membership]\label{lem:c1-actual-membership}
Every exceptional action in Theorem~\ref{thm:three-twentieths}, other
than natural $C_3,A_4$, is accepted by
Lemma~\ref{lem:c1-earlier-fibres}. The action descriptions used there
are certificates on the original permutation sets.
\end{lemma}
\begin{proof}
For the degree-six actions the witnesses are, in order, regular $C_6$,
the degree-six action of $A_4$, $C_3\wr C_2$, and $C_2\times A_4$
with its cyclic three-dimensional binary base. These witnesses also
show that the first and third have normal Sylow $3$-subgroup of index
two, the second is an irreducible coprime affine action outside
degree four, and the fourth has cyclic binary module and elementary
ternary top. Thus they belong to the earlier physical action domains.

The bounded degree-twelve certificate constructs $W=O_2(U)$ and a
Sylow $3$-complement for $12T20,85,164,228,229,265$. It verifies
that $W$ has three actual orbits of size four and local image $V_4$.
Restricting translations to the blocks constructs a literal ambient
$V_4^3$ normalized by $U$. The complement acts faithfully and
transitively on its nine nonidentity local translations, and
transitively on the three blocks. Its order, and the base order, are
\[
\begin{array}{c|rrrrrr}
 U&12T20&12T85&12T164&12T228&12T229&12T265\\\hline
 |W|&4&16&64&64&64&64\\
 |P|&9&9&9&27&27&81.
\end{array}
\]
In the first three rows the complement is $C_3^2$, and an exhibited
base vector has orbit spanning $W$; these are exactly the cyclic
module witnesses used in the proof. In the last three rows the
actual nine-point complement is, respectively, $9T6,9T7,9T17$;
the proper/full binary-base domain applies. In $12T20,12T194$,
the normal ternary subgroup instead has four actual orbits of size
three, local image $C_3$, and full natural $A_4$ block action.
The standard transversal embedding puts the original group inside
$C_3\wr A_4$; the base is diagonal in the first case and full in
the second. Thus $12T20$ also belongs to the earlier prime-base
domain, and $12T194$ belongs to it directly.

These finite statements are part of the bounded action certificate:
they verify actual blocks, local images, full projection and cyclic
generators, rather than deciding membership from order or a catalogue
label. They are separate from the exhaustive normal-pair computation
which produced the high list.  The exact degree-nine order seam used in
Theorem~\ref{thm:three-twentieths} already places its four actions in the
transitive $3$-group owner, and the reversible block-sign dichotomy places
the degree-27 row in that same owner.  Their counting proof therefore
applies uniformly to every actual transitive $3$-group in those degrees;
no further finite action classification is introduced here.
\end{proof}

\begin{theorem}[First split one-triple estimate]\label{thm:c1-first}
The family of split one-triple graphs with exactly one actual regular
$C_3$ orbit and no natural $A_4$ orbit admits a positive forward bound
against complete lower-degree subgroup counts whose row tends to zero
exponentially. After its low cone and the actual orbit families in
Lemma~\ref{lem:c1-earlier-fibres} have been removed, its high cone
is empty.
\end{theorem}
\begin{proof}
Apply Proposition~\ref{prop:c1-low}. For every remaining graph
with an eligible actual orbit, point one such orbit of degree $w$.
Its complete complement $J\leq S_b$, $b=n-w$, includes the external
triple and every other orbit and relation. The exact pointed
comparison is
\[
 \sum_H k(H)=\sum_U\frac{n!}{b!a(U)}
                    \sum_J Z^{\mathrm{surv}}_J(U),
 \qquad 0\leq Z^{\mathrm{surv}}_J(U)\leq Z_J(U).
\]
Here the surviving fibre retains the preceding exclusions on the complete
pointed graph. The unrestricted fibre is used only for the upper bound.
If required, the pattern is retained on each $J$ before using its
rank bound. Deleting any of the indicated nonmarker actions preserves
the unique-$C_3$, no-natural-$A_4$ pattern. Only after the fibre has
been bounded do we enlarge the number of such $J$ to $s_b$.
Normalize once, obtaining terms
\[
 D_U(1+b)\Delta(n,b)2^{\alpha_Ub}\frac{s_b}{L_b}.
\]
For even $w$ the Gaussian coefficient of $b$ is $w/8$; for odd
$w$ its safe value in both parities is $(w-1)/8$. These are strictly
larger than every corresponding $\alpha_U$ in the table. For a check
using only rational numbers, $3^5<2^8$ bounds its five exponents by
$8/15,16/15,25/18,8/9,8/3$, respectively. Each resulting gap is at
least $1/9$. Equation~\eqref{eq:arbitrary-width-pointing} bounds the
coefficient cost for each fixed deletion by a polynomial in $n$.
There are only finitely
many actions in degrees $6,9,12,27$, so the whole forward row is
$O(2^{-\varepsilon n})$ for some $\varepsilon>0$. To form a
disjoint physical partition, assign a graph to the first eligible
family in a fixed order. Summing its intersections before applying
the full pointed upper bound charges each family once. Multiple
eligible orbits give only the positive pointing multiplicity.

Suppose now that a high graph survives. Its residual $K$ has
$20d_3(K)>3m$. Because projection of the complete graph onto $K$
is an isomorphism, the actual nonexternal orbit actions of the graph
are exactly those of $K$. None is natural $C_3$ or $A_4$ by the
pattern, and none is one of the other actions in
Theorem~\ref{thm:three-twentieths}, by
Lemma~\ref{lem:c1-actual-membership} and the preceding removals.
Apply \eqref{eq:orbit-relative-head} to its actual orbit filtration:
\[
 20d_3(K)\leq\sum_i20\delta_{A_i}(N_i)
       \leq\sum_i3\deg A_i=3m,
\]
a contradiction. Fixed points have head zero and keep their degree
in this inequality. Thus every surviving character set
$\Theta(K)$ in the high cone is empty. This is a statement about
complete marked graphs; it introduces no additional numerical
kernel or scalar term.
\end{proof}

\paragraph{The two exact incidence tests.}
The preceding proof respects Proposition~\ref{prop:c1-incidence} on
every disjoint piece. In particular, applying that proposition with
$\chi\in\Theta(K)$ gives
\[
 \sum_K\sigma^{\mathrm{surv}}_3(K)
 =\sum_{|x|=3}\ 
 \sum_{\substack{N^x=N,\ N\cap\langle x\rangle=1\\
        \chi_{N,x}\in\Theta(N\langle x\rangle)}}
 \frac1{|\{y\in Nx:|y|=3\}|}.
\]
The reciprocal is preserved. It is not harmless to drop it: for
$K_0=C_2^{3h}\rtimes\langle x\rangle$, with $x$ cycling $h$
triples of pair coordinates, and $K=K_0\times A_5^2$ on
$m=6h+10$ points, its indicated quotient coset contains exactly
$4^h\cdot21^2$ elements of order three. This already loses
$m/3-O(1)$ bits if the reciprocal is discarded.

Likewise the affine test on the same source
$J=\mathbb F_4^r\rtimes C_3$ has, above the fixed top map to
$\mathbb F_4^t\rtimes C_3$, exactly
$4^{t(r+1)}$ cocycles and
$4^t\prod_{i=0}^{t-1}(4^r-4^i)$ onto lifts. A second mark stays
on this same $J$. In the proof above the external triple is part of
$J$ in every earlier fibre, including the degree-six affine fibre;
it is never appended as an independent character maximum. The final
zero-set argument makes all these joint high surviving incidences
empty and therefore requires no factorization of their moments.
